% METADADOS DO ARTIGO
% ===================
% Este é o único arquivo em que autores preenchem dados editoriais. Substitua
% todos os campos por informações verificadas. Os dados abaixo são FICTÍCIOS e
% existem apenas para que este modelo possa ser visualizado antes do envio.

\newcommand{\ArtigoTitulo}{Lorem ipsum: demonstração de estrutura para artigos científicos}
\newcommand{\ArtigoSubtitulo}{Texto ilustrativo para revisão editorial e diagramação}

% Liste autores na ordem de autoria e use sobrescritos para associá-los às
% filiações. Mantenha asterisco somente para a pessoa correspondente.
\newcommand{\ArtigoAutores}{Autora de Exemplo\textsuperscript{1*}; Coautor de Exemplo\textsuperscript{2}}
\newcommand{\ArtigoAfiliacoes}{\textsuperscript{1}Instituição de Exemplo, Campus Demonstrativo, Cidade, BA, Brasil.\\\textsuperscript{2}Instituição de Exemplo, Campus Ilustrativo, Cidade, BA, Brasil.}
\newcommand{\ArtigoContatos}{E-mail: autora.exemplo@exemplo.org. ORCID: https://orcid.org/0000-0000-0000-0000.}

% Resumo em português: máximo de 250 palavras; deve informar objetivo e os
% principais resultados. Palavras-chave: até cinco, sem repetir o título.
\newcommand{\ArtigoResumo}{Este texto fictício demonstra a organização prevista para a submissão de artigos no livro. Seu objetivo é permitir que autores e equipe de diagramação observem a abertura, as seções, as citações, as figuras e as referências antes da substituição pelo manuscrito definitivo. Não apresenta resultados de pesquisa nem informações institucionais reais.}
\newcommand{\ArtigoPalavrasChave}{Diagramação. Tipografia. Manuscrito. Demonstração.}

% Traduções fiéis do título, resumo e palavras-chave em português.
\newcommand{\ArtigoTituloIngles}{Lorem ipsum: a structural demonstration for scientific articles}
\newcommand{\ArtigoAbstract}{This fictional text demonstrates the organisation required for articles in the book. It lets authors and the layout team inspect the opening matter, sections, citations, figures, and references before replacing it with the final manuscript. It reports neither research results nor real institutional information.}
\newcommand{\ArtigoKeywords}{Typesetting. Typography. Manuscript. Demonstration.}
